\PassOptionsToPackage{numbers}{natbib}
\documentclass{article}
\usepackage{iclr2024_conference}
\usepackage{luatexja-fontspec}
\usepackage{hyperref}
\usepackage{url}
\usepackage{booktabs}
\usepackage{amsmath}
\usepackage{amsfonts}
\usepackage{graphicx}
\graphicspath{{../}}

\InputIfFileExists{values.tex}{}{}
\providecommand{\airasval}[1]{\textbf{??airasval:\detokenize{#1}??}}
\providecommand{\unverified}[1]{#1}
\providecommand{\airasrecordlink}[1]{#1}

\title{SciGym における反復上限の効果の事前登録評価: open-weight 3 モデルで上限を 1 から 80 まで振り、改善と飽和を測る}
\author{AIRAS}

\begin{document}
\maketitle

\begin{abstract}
SciGym \citep{duan-2025-measuring} は、反応をすべて取り除いた SBML モデルを LLM エージェントに渡し、摂動実験とコード実行を繰り返して反応機構を復元させるベンチマークである。論文はエージェントが使える反復の上限を 20 に固定して評価しており、上限を変えたときに性能がどう変わるかは報告されていない。エージェントは残り反復数を毎回知らされるので、上限は探索の深さそのものを決めるパラメータである。本研究は、公開コードの Controller と評価器をそのまま用い、反復上限だけを 1 / 5 / 10 / 20 / 40 / 80 に変えて、open-weight 3 モデル（DeepSeek-V4.1-Flash、Qwen3.8-27B、Nemotron-3-Ultra）を RIKYU の GB200 ノード上で自前 serving し、SciGym-small 137 件と SciGym-large 213 件で評価する。実験前に、上限を論文の 20 から 80 に増やすと平均 Simulation Trajectory Error（STE）が雑音幅を超えて改善すること（仮説 1）と、40 から 80 への改善は 0.05 未満にとどまること（仮説 2）を、モデル × split ごとの 12 の主張として登録する。上限に対して実際に使った反復数、自発的に提出した件数、文脈長の使用量と超過件数を STE と並べて報告し、上限が効くのはモデルが予算を使い切る場合に限られるかを見る。
\end{abstract}

\section{はじめに}
\label{sec:intro}
LLM に科学的発見のサイクル全体（実験の設計、結果の解析、機構の提案）を行わせて測るベンチマークとして、SciGym \citep{duan-2025-measuring} が提案された。BioModels 由来の SBML モデルから反応をすべて取り除いたものをエージェントに渡し、エージェントは初期濃度を変えた摂動実験を要求し、返ってきた時系列を Python で解析し、最終的な SBML を提出する。提出モデルは真のモデルと、軌道誤差（STE）と反応の一致度（RMS）で比較される。

論文の主結果である Table 1 は、3 社の pro / mini 計 6 モデルを SciGym-small の 137 件で評価した平均値である \cite[s1.p2, s1.p6]{duan-2025-measuring}。この評価では、エージェントはいつでも提出できるが反復の上限は 20 に固定されている \cite[s1.p3]{duan-2025-measuring}。AIRAS の先行研究は、RIKYU の推論 API で提供される open-weight モデル 6 種を、公開コードの Controller で SciGym-small 137 件を最大 20 反復で解かせて評価し \cite[s3.p1]{scigymrikyuopenmodels-2026-scigym}、deepseek-v4.1-flash の平均 STE が \unverified{0.2405} \cite[s3.p2]{scigymrikyuopenmodels-2026-scigym}、qwen3.8-27b が \unverified{0.2479} \cite[s3.p3]{scigymrikyuopenmodels-2026-scigym} と、論文 Table 1 の最良行 Gemini-2.5-Pro の \unverified{0.3212} \cite[s1.p5]{duan-2025-measuring} を下回ることを示した。

いずれの評価も反復上限 20 という 1 点での測定である。SciGym の Controller はエージェントに毎ターン「残り反復数 / 上限」を伝えるので、上限は単に打ち切りの位置ではなく、エージェントがどれだけ探索してから提出するかを決める条件でもある。上限を増やせば実験データが増えて性能が上がるのか、それともモデルは上限に関わらず一定の反復で自発的に提出してしまい上限は効かないのか、そして効くとすればどこで飽和するのかは分かっていない。本研究はこの 1 軸を系統的に振る。判定は AIRAS の事前登録の枠組みに従い、仮説・判定基準・予測区間・実行条件を実験前に記録へ凍結し、実験後は計測値だけを流し込む。

\section{関連研究}
\label{sec:related}
SciGym \citep{duan-2025-measuring} は、単一スキルを測る既存のベンチマークと異なり、エージェントが自ら摂動実験を設計してその結果を受け取る閉ループを評価する。評価指標は STE、RMS（modifier まで要求する厳格版もある）、NTS の 3 つで、論文は pro モデルが mini モデルを一貫して上回り、Gemini-Pro が最良であったと報告している \cite[s1.p6]{duan-2025-measuring}。公開コードの既定の config は反復上限 5 \cite[s2.p2]{h4duan-2025-scigym}、デバッグ反復 5 \cite[s2.p3]{h4duan-2025-scigym} で、論文の条件（20 と 3）とは異なる。

open-weight モデルの 20 反復での評価では、deepseek-v4.1-flash が \unverified{0.2405} \cite[s3.p2]{scigymrikyuopenmodels-2026-scigym}、qwen3.8-27b が \unverified{0.2479} \cite[s3.p3]{scigymrikyuopenmodels-2026-scigym} に達している。本研究の 3 モデルのうち 2 つはこの先行研究と同じモデルであり、上限 20 の run は先行値の再測定にもなる。

\section{仮説と予測}
\label{sec:hypothesis}
以下は記録（record.json）から生成された仮説・主張・判定基準・予測区間の一覧である。実験後は Observed と Verdict が計測値から埋められる。

% AUTO-GENERATED by the update_record tool. DO NOT EDIT.
% verify_paper regenerates this file from record.json and the run outputs
% and fails on any difference, so manual edits always surface.
\noindent\textbf{Hypothesis 1.} SciGym 公開コードの Controller を論文条件（eval\_debug\_rounds=3、temperature=1.0）のまま用い、open-weight 3 モデル（DeepSeek-V4.1-Flash、Qwen3.8-27B、Nemotron-3-Ultra）を自前 serving して SciGym-small 137 件と SciGym-large 213 件を解かせたとき、反復上限を論文の 20 から 80 に増やすと平均 Simulation Trajectory Error（STE）は雑音幅を超えて（0.03 以上）改善する。~\cite[s1.p3, s1.p5, s1.p6]{duan-2025-measuring}, \cite[s3.p1]{scigymrikyuopenmodels-2026-scigym}
\begin{enumerate}
\item[\textbf{Claim 1}] DeepSeek-V4.1-Flash で SciGym-small 137 件を解かせたとき、反復上限 80 の平均 STE は反復上限 20（論文条件）より 0.03 以上低い。~\cite[s1.p3, s1.p4, s1.p5]{duan-2025-measuring}, \cite[s2.p1]{h4duan-2025-scigym}
  \emph{Rationale:} DeepSeek-V4.1-Flash の small split で、論文条件を超える予算が雑音幅（137 件平均の標準誤差約 0.02、差で約 0.03）を超えて効くことを示し、h1 の当該モデル・split の部分を担う。
  \emph{Criterion:} \texttt{\detokenize{deepseek-v4.1-flash-small-i80}}.\texttt{\detokenize{trajectory_smape}} $-$ \texttt{\detokenize{deepseek-v4.1-flash-small-i20}}.\texttt{\detokenize{trajectory_smape}} $\leq -0.03$.
  \emph{Prediction:} $[-0.15, -0.02]$ (先行する 100 反復の内部試行では kimi-k3 が 20$\to$100 反復で STE 0.225$\to$0.088（平均 67.8 反復を使用）、qwen3.6-35b は平均 20.7 反復で自発提出し 0.457$\to$0.446。予算を使い切るモデルほど改善が大きいと見て、モデルごとに幅を置く).
  \emph{Observed:} -0.0643907.
  \emph{Verdict:} supported.
\item[\textbf{Claim 2}] DeepSeek-V4.1-Flash で SciGym-large 213 件を解かせたとき、反復上限 80 の平均 STE は反復上限 20（論文条件）より 0.03 以上低い。~\cite[s1.p3, s1.p4, s1.p5]{duan-2025-measuring}, \cite[s2.p1]{h4duan-2025-scigym}
  \emph{Rationale:} DeepSeek-V4.1-Flash の large split で、論文条件を超える予算が雑音幅（137 件平均の標準誤差約 0.02、差で約 0.03）を超えて効くことを示し、h1 の当該モデル・split の部分を担う。
  \emph{Criterion:} \texttt{\detokenize{deepseek-v4.1-flash-large-i80}}.\texttt{\detokenize{trajectory_smape}} $-$ \texttt{\detokenize{deepseek-v4.1-flash-large-i20}}.\texttt{\detokenize{trajectory_smape}} $\leq -0.03$.
  \emph{Prediction:} $[-0.15, -0.02]$ (先行する 100 反復の内部試行では kimi-k3 が 20$\to$100 反復で STE 0.225$\to$0.088（平均 67.8 反復を使用）、qwen3.6-35b は平均 20.7 反復で自発提出し 0.457$\to$0.446。予算を使い切るモデルほど改善が大きいと見て、モデルごとに幅を置く).
  \emph{Observed:} -0.1495.
  \emph{Verdict:} supported.
\item[\textbf{Claim 3}] Qwen3.8-27B で SciGym-small 137 件を解かせたとき、反復上限 80 の平均 STE は反復上限 20（論文条件）より 0.03 以上低い。~\cite[s1.p3, s1.p4, s1.p5]{duan-2025-measuring}, \cite[s2.p1]{h4duan-2025-scigym}
  \emph{Rationale:} Qwen3.8-27B の small split で、論文条件を超える予算が雑音幅（137 件平均の標準誤差約 0.02、差で約 0.03）を超えて効くことを示し、h1 の当該モデル・split の部分を担う。
  \emph{Criterion:} \texttt{\detokenize{qwen3.8-27b-small-i80}}.\texttt{\detokenize{trajectory_smape}} $-$ \texttt{\detokenize{qwen3.8-27b-small-i20}}.\texttt{\detokenize{trajectory_smape}} $\leq -0.03$.
  \emph{Prediction:} $[-0.08, 0.02]$ (先行する 100 反復の内部試行では kimi-k3 が 20$\to$100 反復で STE 0.225$\to$0.088（平均 67.8 反復を使用）、qwen3.6-35b は平均 20.7 反復で自発提出し 0.457$\to$0.446。予算を使い切るモデルほど改善が大きいと見て、モデルごとに幅を置く).
  \emph{Observed:} -0.0973588.
  \emph{Verdict:} supported.
\item[\textbf{Claim 4}] Qwen3.8-27B で SciGym-large 213 件を解かせたとき、反復上限 80 の平均 STE は反復上限 20（論文条件）より 0.03 以上低い。~\cite[s1.p3, s1.p4, s1.p5]{duan-2025-measuring}, \cite[s2.p1]{h4duan-2025-scigym}
  \emph{Rationale:} Qwen3.8-27B の large split で、論文条件を超える予算が雑音幅（137 件平均の標準誤差約 0.02、差で約 0.03）を超えて効くことを示し、h1 の当該モデル・split の部分を担う。
  \emph{Criterion:} \texttt{\detokenize{qwen3.8-27b-large-i80}}.\texttt{\detokenize{trajectory_smape}} $-$ \texttt{\detokenize{qwen3.8-27b-large-i20}}.\texttt{\detokenize{trajectory_smape}} $\leq -0.03$.
  \emph{Prediction:} $[-0.08, 0.02]$ (先行する 100 反復の内部試行では kimi-k3 が 20$\to$100 反復で STE 0.225$\to$0.088（平均 67.8 反復を使用）、qwen3.6-35b は平均 20.7 反復で自発提出し 0.457$\to$0.446。予算を使い切るモデルほど改善が大きいと見て、モデルごとに幅を置く).
  \emph{Observed:} -0.130822.
  \emph{Verdict:} supported.
\item[\textbf{Claim 5}] Nemotron-3-Ultra で SciGym-small 137 件を解かせたとき、反復上限 80 の平均 STE は反復上限 20（論文条件）より 0.03 以上低い。~\cite[s1.p3, s1.p4, s1.p5]{duan-2025-measuring}, \cite[s2.p1]{h4duan-2025-scigym}
  \emph{Rationale:} Nemotron-3-Ultra の small split で、論文条件を超える予算が雑音幅（137 件平均の標準誤差約 0.02、差で約 0.03）を超えて効くことを示し、h1 の当該モデル・split の部分を担う。
  \emph{Criterion:} \texttt{\detokenize{nemotron-3-ultra-small-i80}}.\texttt{\detokenize{trajectory_smape}} $-$ \texttt{\detokenize{nemotron-3-ultra-small-i20}}.\texttt{\detokenize{trajectory_smape}} $\leq -0.03$.
  \emph{Prediction:} $[-0.12, -0.02]$ (先行する 100 反復の内部試行では kimi-k3 が 20$\to$100 反復で STE 0.225$\to$0.088（平均 67.8 反復を使用）、qwen3.6-35b は平均 20.7 反復で自発提出し 0.457$\to$0.446。予算を使い切るモデルほど改善が大きいと見て、モデルごとに幅を置く).
  \emph{Observed:} -0.0280823.
  \emph{Verdict:} refuted.
\item[\textbf{Claim 6}] Nemotron-3-Ultra で SciGym-large 213 件を解かせたとき、反復上限 80 の平均 STE は反復上限 20（論文条件）より 0.03 以上低い。~\cite[s1.p3, s1.p4, s1.p5]{duan-2025-measuring}, \cite[s2.p1]{h4duan-2025-scigym}
  \emph{Rationale:} Nemotron-3-Ultra の large split で、論文条件を超える予算が雑音幅（137 件平均の標準誤差約 0.02、差で約 0.03）を超えて効くことを示し、h1 の当該モデル・split の部分を担う。
  \emph{Criterion:} \texttt{\detokenize{nemotron-3-ultra-large-i80}}.\texttt{\detokenize{trajectory_smape}} $-$ \texttt{\detokenize{nemotron-3-ultra-large-i20}}.\texttt{\detokenize{trajectory_smape}} $\leq -0.03$.
  \emph{Prediction:} $[-0.12, -0.02]$ (先行する 100 反復の内部試行では kimi-k3 が 20$\to$100 反復で STE 0.225$\to$0.088（平均 67.8 反復を使用）、qwen3.6-35b は平均 20.7 反復で自発提出し 0.457$\to$0.446。予算を使い切るモデルほど改善が大きいと見て、モデルごとに幅を置く).
  \emph{Observed:} -0.0411857.
  \emph{Verdict:} supported.
\end{enumerate}
\noindent\emph{Assumed, so that the claims together imply Hypothesis 1:}
\begin{itemize}
\item temperature 1.0 での 1 回実行の平均 STE の実行間変動が判定に対して十分小さいこと（137 件平均の標準誤差約 0.02）（c1〜c12）
\item 自前 serving（vLLM / NIM）したモデルが、提供元の推奨設定で呼んだ場合と同じ能力で応答すること。思考はサーバー側の reasoning parser で本文から分離する（c1〜c12）
\item 文脈長の拡張（YaRN、NIM\_MAX\_MODEL\_LEN）が 262,144 以内の応答品質を変えないこと（c1〜c12）
\item aarch64 向けの libroadrunner 2.7.0 と antimony 2.14.0 が、論文環境と同じ軌道と評価値を与えること（c1〜c12）
\end{itemize}
\noindent\textbf{Hypothesis 2.} 同じ条件で反復上限を 40 から 80 に倍増しても、平均 STE の改善は 0.05 未満にとどまる（論文条件の 2 倍程度で性能は飽和する）。~\cite[s1.p3]{duan-2025-measuring}
\begin{enumerate}
\item[\textbf{Claim 7}] DeepSeek-V4.1-Flash で SciGym-small 137 件を解かせたとき、反復上限 40 から 80 への平均 STE の改善は 0.05 未満である。~\cite[s1.p3, s1.p4]{duan-2025-measuring}, \cite[s2.p1]{h4duan-2025-scigym}
  \emph{Rationale:} DeepSeek-V4.1-Flash の small split で、予算を 40 から 80 に倍増しても伸びが判定幅に収まる（飽和している）ことを示し、h2 の当該モデル・split の部分を担う。基準 run は c1 の反復上限 80。
  \emph{Criterion:} \texttt{\detokenize{deepseek-v4.1-flash-small-i40}}.\texttt{\detokenize{trajectory_smape}} $-$ \texttt{\detokenize{deepseek-v4.1-flash-small-i80}}.\texttt{\detokenize{trajectory_smape}} $\leq 0.05$.
  \emph{Prediction:} $[0, 0.08]$ (内部試行では kimi-k3 が 100 反復でも平均 67.8 反復で提出しており、40 を超える予算の使用は部分的。qwen 系は 20 前後で自発提出する。40$\to$80 の差は雑音幅（約 0.03）程度と見る).
  \emph{Observed:} 0.0390427.
  \emph{Verdict:} supported.
\item[\textbf{Claim 8}] DeepSeek-V4.1-Flash で SciGym-large 213 件を解かせたとき、反復上限 40 から 80 への平均 STE の改善は 0.05 未満である。~\cite[s1.p3, s1.p4]{duan-2025-measuring}, \cite[s2.p1]{h4duan-2025-scigym}
  \emph{Rationale:} DeepSeek-V4.1-Flash の large split で、予算を 40 から 80 に倍増しても伸びが判定幅に収まる（飽和している）ことを示し、h2 の当該モデル・split の部分を担う。基準 run は c2 の反復上限 80。
  \emph{Criterion:} \texttt{\detokenize{deepseek-v4.1-flash-large-i40}}.\texttt{\detokenize{trajectory_smape}} $-$ \texttt{\detokenize{deepseek-v4.1-flash-large-i80}}.\texttt{\detokenize{trajectory_smape}} $\leq 0.05$.
  \emph{Prediction:} $[0, 0.08]$ (内部試行では kimi-k3 が 100 反復でも平均 67.8 反復で提出しており、40 を超える予算の使用は部分的。qwen 系は 20 前後で自発提出する。40$\to$80 の差は雑音幅（約 0.03）程度と見る).
  \emph{Observed:} 0.0583697.
  \emph{Verdict:} refuted.
\item[\textbf{Claim 9}] Qwen3.8-27B で SciGym-small 137 件を解かせたとき、反復上限 40 から 80 への平均 STE の改善は 0.05 未満である。~\cite[s1.p3, s1.p4]{duan-2025-measuring}, \cite[s2.p1]{h4duan-2025-scigym}
  \emph{Rationale:} Qwen3.8-27B の small split で、予算を 40 から 80 に倍増しても伸びが判定幅に収まる（飽和している）ことを示し、h2 の当該モデル・split の部分を担う。基準 run は c3 の反復上限 80。
  \emph{Criterion:} \texttt{\detokenize{qwen3.8-27b-small-i40}}.\texttt{\detokenize{trajectory_smape}} $-$ \texttt{\detokenize{qwen3.8-27b-small-i80}}.\texttt{\detokenize{trajectory_smape}} $\leq 0.05$.
  \emph{Prediction:} $[-0.02, 0.04]$ (内部試行では kimi-k3 が 100 反復でも平均 67.8 反復で提出しており、40 を超える予算の使用は部分的。qwen 系は 20 前後で自発提出する。40$\to$80 の差は雑音幅（約 0.03）程度と見る).
  \emph{Observed:} 0.0443245.
  \emph{Verdict:} supported.
\item[\textbf{Claim 10}] Qwen3.8-27B で SciGym-large 213 件を解かせたとき、反復上限 40 から 80 への平均 STE の改善は 0.05 未満である。~\cite[s1.p3, s1.p4]{duan-2025-measuring}, \cite[s2.p1]{h4duan-2025-scigym}
  \emph{Rationale:} Qwen3.8-27B の large split で、予算を 40 から 80 に倍増しても伸びが判定幅に収まる（飽和している）ことを示し、h2 の当該モデル・split の部分を担う。基準 run は c4 の反復上限 80。
  \emph{Criterion:} \texttt{\detokenize{qwen3.8-27b-large-i40}}.\texttt{\detokenize{trajectory_smape}} $-$ \texttt{\detokenize{qwen3.8-27b-large-i80}}.\texttt{\detokenize{trajectory_smape}} $\leq 0.05$.
  \emph{Prediction:} $[-0.02, 0.04]$ (内部試行では kimi-k3 が 100 反復でも平均 67.8 反復で提出しており、40 を超える予算の使用は部分的。qwen 系は 20 前後で自発提出する。40$\to$80 の差は雑音幅（約 0.03）程度と見る).
  \emph{Observed:} 0.07322.
  \emph{Verdict:} refuted.
\item[\textbf{Claim 11}] Nemotron-3-Ultra で SciGym-small 137 件を解かせたとき、反復上限 40 から 80 への平均 STE の改善は 0.05 未満である。~\cite[s1.p3, s1.p4]{duan-2025-measuring}, \cite[s2.p1]{h4duan-2025-scigym}
  \emph{Rationale:} Nemotron-3-Ultra の small split で、予算を 40 から 80 に倍増しても伸びが判定幅に収まる（飽和している）ことを示し、h2 の当該モデル・split の部分を担う。基準 run は c5 の反復上限 80。
  \emph{Criterion:} \texttt{\detokenize{nemotron-3-ultra-small-i40}}.\texttt{\detokenize{trajectory_smape}} $-$ \texttt{\detokenize{nemotron-3-ultra-small-i80}}.\texttt{\detokenize{trajectory_smape}} $\leq 0.05$.
  \emph{Prediction:} $[-0.01, 0.06]$ (内部試行では kimi-k3 が 100 反復でも平均 67.8 反復で提出しており、40 を超える予算の使用は部分的。qwen 系は 20 前後で自発提出する。40$\to$80 の差は雑音幅（約 0.03）程度と見る).
  \emph{Observed:} 0.0264056.
  \emph{Verdict:} supported.
\item[\textbf{Claim 12}] Nemotron-3-Ultra で SciGym-large 213 件を解かせたとき、反復上限 40 から 80 への平均 STE の改善は 0.05 未満である。~\cite[s1.p3, s1.p4]{duan-2025-measuring}, \cite[s2.p1]{h4duan-2025-scigym}
  \emph{Rationale:} Nemotron-3-Ultra の large split で、予算を 40 から 80 に倍増しても伸びが判定幅に収まる（飽和している）ことを示し、h2 の当該モデル・split の部分を担う。基準 run は c6 の反復上限 80。
  \emph{Criterion:} \texttt{\detokenize{nemotron-3-ultra-large-i40}}.\texttt{\detokenize{trajectory_smape}} $-$ \texttt{\detokenize{nemotron-3-ultra-large-i80}}.\texttt{\detokenize{trajectory_smape}} $\leq 0.05$.
  \emph{Prediction:} $[-0.01, 0.06]$ (内部試行では kimi-k3 が 100 反復でも平均 67.8 反復で提出しており、40 を超える予算の使用は部分的。qwen 系は 20 前後で自発提出する。40$\to$80 の差は雑音幅（約 0.03）程度と見る).
  \emph{Observed:} 0.0152157.
  \emph{Verdict:} supported.
\end{enumerate}
\noindent\emph{Assumed, so that the claims together imply Hypothesis 2:}
\begin{itemize}
\item h1 と同じ実行・serving・評価環境の仮定（c7〜c12）
\item 反復上限 40 と 80 の run が同じサーバー設定・同じコミットで実行され、差が予算以外の条件によらないこと（c7〜c12）
\end{itemize}


\paragraph{判定基準と予測区間の根拠.}
仮説 1（Claim 1〜6）の反証線は「（上限 80 の平均 STE）$-$（上限 20 の平均 STE）$\le -0.03$」である。先行する Table 1 再現では件ごとの STE の標準偏差はおよそ 0.27 で、137 件平均の標準誤差は約 0.02、2 つの run の差では約 0.03 になる。margin $-0.03$ は、この雑音幅を超える改善だけを「効いた」と数えるための線である。仮説 2（Claim 7〜12）の反証線は「（上限 40 の平均 STE）$-$（上限 80 の平均 STE）$\le 0.05$」で、予算を倍にしても改善が 0.05 未満であることを飽和の定義とする。上限 80 の run は仮説 1 の主張の下にあり、仮説 2 の主張はそれを基準 run として参照する。large（213 件）は標準誤差がやや小さいが、同じ margin を使う。

予測区間はモデルごとに置いた。根拠は、本記録には含まれない予備試行で、kimi-k3 が上限 100 のとき平均 67.8 反復まで使って STE が 20 反復時の 0.225 から 0.088 へ大きく改善した一方、qwen3.6-35b は上限 100 でも平均 20.7 反復で自発的に提出し 0.457 から 0.446 とほぼ横ばいだったことである。予算を使い切るモデルほど改善が大きいと見て、DeepSeek-V4.1-Flash（文脈長 1M、大型 MoE）には広い改善幅、Qwen3.8-27B（qwen 系は早く提出する傾向）には 0 を含む幅、Nemotron-3-Ultra（20 反復の先行実行で 137 件中 50 件を上限前に提出）にはその中間を置く。反復上限 1 / 5 / 10 の run は主張の判定には使わず、曲線の低予算側を表と図で示すためのものである。

\section{方法}
\label{sec:method}
\paragraph{タスク.}
各インスタンスは真の SBML（\texttt{truth.xml}）、反応を取り除いた不完全な SBML（\texttt{partial.xml}）、シミュレーション条件（\texttt{truth.sedml}）からなる。エージェントは不完全な SBML を受け取り、各ターンで実験（初期濃度の変更を指定し、真のモデルの時系列を得る）、コード実行（numpy / pandas / scipy / libsbml と、仮説モデルを走らせる \texttt{simulate} 関数が使える）、提出のいずれかを markdown 形式で返す。提出は任意の時点で行え、論文では反復数の上限は 20 である \cite[s1.p3]{duan-2025-measuring}。提出した SBML が無効な場合は追加のデバッグ反復が許され \cite[s1.p4]{duan-2025-measuring}、それでも無効なら不完全 SBML のまま採点される。

\paragraph{反復上限の意味.}
公開コードの Controller は、上限 $k$ のとき反復 $0$ から $k$ までの $k+1$ 回の実験・コード実行を許し、各ターンのフィードバックに「残り反復数 / 上限」を含める。残りが 0 になると最終モデルの提出を促す。上限はプロンプトの一部としてエージェントに見えているので、上限 80 の run の 20 反復目と上限 20 の run の 20 反復目は同じ状態ではない。したがって 1 本の長い run から途中の点を切り出すことはできず、上限ごとに独立した run を実行する。

\paragraph{実行系.}
ループ、プロンプト、実験の実行はすべて公開コード \citep{h4duan-2025-scigym} の \texttt{Controller} をそのまま用いる。本研究が書くのは、(1) OpenAI 互換 API でモデルを呼ぶ \texttt{LLM} 実装（公式の GPT 実装と同じく、システムプロンプトと全履歴を毎回送る）、(2) インスタンスごとに 1 プロセスで Controller を起動し、複数件を並列に走らせる駆動部、(3) 各インスタンスの提出モデルと公式データを評価層の入力ファイルに書き出す部分だけである。公式コードでは履歴がクラス属性として定義されており 1 プロセスで複数件を走らせると履歴が持ち越されるが、1 件 1 プロセスとすることでこの問題を避ける。LLM が書いた解析コードや提出モデルの評価が ODE ソルバーの C ライブラリ内で止まると公式の 3 分制限（\texttt{signal.alarm}）が効かないので、駆動部は出力が 60 分止まった試行を打ち切って反復 0 からやり直す（最大 3 試行）。1 件 1 試行の上限時間は上限反復に比例させ、$\max(2\,\text{時間},\ 12\,\text{分} \times \text{上限})$ とする。

\paragraph{serving.}
3 モデルはいずれも RIKYU の GB200 ノード（4 GPU）上に、run ごとに専用の OpenAI 互換サーバーを Slurm job で立てて呼ぶ。DeepSeek-V4.1-Flash と Qwen3.8-27B は vLLM 0.30.0（tensor parallel 4）、Nemotron-3-Ultra は NVIDIA NIM 2.0.11（vLLM backend）である。重みは共有ストレージに置かず、job の開始時に Hugging Face から計算ノードのローカルディスクへ取得する。思考（reasoning）はサーバー側の reasoning parser で本文から分離し（vLLM は \texttt{qwen3} / \texttt{deepseek\_v41}、NIM は既定）、分離し損ねて本文の先頭に思考の節が残った応答は駆動部が除いてから Controller に渡し、その件数を表に出す。サーバーの job は計算機の 4 日上限で打ち切られうるので、実験側はサーバーの接続先を共有ディスク上のファイルから読み、接続断のたびに読み直して別ノードで立ち直ったサーバーに繋ぎ直す。

\paragraph{文脈長.}
反復が増えると会話履歴に実験データの表が積み上がり、上限 80 の large では配布時の文脈長 262,144 トークンを超える件が多くなると見込まれる。文脈超過が反復の効果と混ざるのを避けるため、論文条件からの逸脱として文脈長を 1M に揃える。DeepSeek-V4.1-Flash は native で 1,048,576 トークン、Qwen3.8-27B は公式モデルカードの YaRN 設定（factor 4.0、original 262,144）で 1,000,000 トークン、Nemotron-3-Ultra は NIM の \texttt{NIM\_MAX\_MODEL\_LEN} で 1,048,576 トークンに設定する（後の 2 つは配布時の既定 262,144 からの拡張）。それでも超過した件は提出なしとして不完全モデルの採点値で平均に含め、件数と 1 呼び出しあたりの平均・最大プロンプト長を表に出す。

\paragraph{指標.}
採点は評価層 airas-eval のベンチマークパック \texttt{scigym\_small}（137 件）と \texttt{scigym\_large}（213 件）が行う。パックは公式リリースの各件（truth / partial / sedml）を sha256 で固定し、公開コードの \texttt{Evaluator}（コミット 88a7b93）で提出モデルを採点して、STE（\texttt{trajectory\_smape}）、RMS の適合率・再現率・F1（modifier なし \texttt{reaction\_*}、あり \texttt{reaction\_*\_with\_modifiers}）、NTS の F1（\texttt{topology\_f1}）を件の単純平均で返す。提出 SBML が無い、または無効な件は、公式評価どおり不完全 SBML の採点値で平均に含める。駆動部は STE に加えて、提出までに使った反復数の平均、上限前に自ら提出した件数、LLM 呼び出し数、入力・出力トークン、1 呼び出しあたりの平均・最大プロンプト長、文脈超過の件数を数える。

\section{実験設計}
\label{sec:design}
\paragraph{独立変数.}
反復上限 \texttt{max\_iterations} $\in \{1, 5, 10, 20, 40, 80\}$。20 が論文条件 \cite[s1.p3]{duan-2025-measuring} で、それ以外は 2 倍刻みの前後である。上限 1 は観測 2 回で提出する下限、80 は文脈長 1M と 4 日の job 上限の中で収まる上限として置いた。

\paragraph{モデル.}
DeepSeek-V4.1-Flash（deepseek-ai/DeepSeek-V4.1-Flash、FP8、native 1M 文脈）、Qwen3.8-27B（Qwen/Qwen3.8-27B、BF16）、Nemotron-3-Ultra（nvidia/nemotron-3-ultra-550b-a55b、NIM の NVFP4 量子化）の 3 種。いずれも重みが公開されており、思考は既定の有効のまま、温度は公開 config の 1.0 \cite[s2.p1]{h4duan-2025-scigym} で呼ぶ。

\paragraph{データ.}
SciGym の公式リリース（Hugging Face \texttt{h4duan/scigym-sbml}）のうち、論文の評価対象は small split の 137 件である \cite[s1.p2]{duan-2025-measuring}。本研究はこれに加えて公式リリースの large split 213 件も用いる（いずれも RIKYU 上のコピー）。評価層は各件の内容を公式リリースの sha256 と照合する。

\paragraph{run.}
1 run は 1 モデル × 1 split × 1 上限で、3 × 2 × 6 = 36 run。run id は \texttt{<model>-<split>-i<上限 2 桁>}（例 \texttt{qwen3.8-27b-small-i20}）。各 run は各インスタンスを 1 回ずつ実行する。並列は 1 run あたり 16 件で、run ごとに専用のサーバー 1 ノードを使う。

\paragraph{設定.}
\texttt{eval\_debug\_rounds} は論文本文の 3 \cite[s1.p4]{duan-2025-measuring}（公開 config の既定値は 5 \cite[s2.p3]{h4duan-2025-scigym}）、\texttt{temperature} は 1.0 \cite[s2.p1]{h4duan-2025-scigym}、摂動は初期濃度の変更のみとする。応答のトークン上限は 32768 とし、思考で使い切った応答は文脈長に収まる範囲で上限を倍にして呼び直す。上限反復以外の設定は全 run で同一である。

\paragraph{計算資源と環境.}
実行は Seyval を通じて RIKYU（Slurm、GB200 ノード、aarch64）で行い、LLM は同じ計算機上のサーバーを呼ぶ。環境は Dockerfile と \texttt{uv.lock} で固定する。aarch64 向けの wheel が無いため、libroadrunner は公式の 2.8.0 ではなく 2.7.0、antimony は 2.14.0 を用い、rrplugins と phrasedml は除外する。この差でシミュレーション結果がわずかに変わる可能性があり、仮説の前提として記録する。

\paragraph{段階と順序.}
sanity は small の先頭 1 件を 2 反復で走らせて配管を確かめる。主張の判定は full の結果だけで行い、sanity の結果は記録に取り込まない。full は上限 1 の 6 run を先に回し、終わり次第残り 30 run を同時に投入する（実行基盤の同時 run 上限が 30）。

\paragraph{run と主張の対応.}
Claim 1〜6 は model × split（DeepSeek small / large、Qwen small / large、Nemotron small / large の順）ごとに上限 1 / 5 / 10 / 20 / 80 の 5 run を持ち、判定は上限 80 と 20 の差で行う。Claim 7〜12 は同じ順で上限 40 の run を 1 つずつ持ち、判定は対応する Claim の上限 80 の run との差で行う。

% ===== airas: post-experiment sections, filled once runs exist =====

\section{結果（暫定版 v2）}
\label{sec:results}
本版は 2026 年 10 月 2 日（JST）時点の暫定版 v2 である。36 run のうち 32 run の結果を取り込み、残る 4 run（DeepSeek-V4.1-Flash と Qwen3.8-27B の上限 80、small と large）は同じサーバー群で実行中である。Nemotron-3-Ultra は 12 run がすべて揃い、Claim 5・6・11・12 の判定が記録から導かれた。他の 8 主張は上限 80 の run を待って判定保留である。表 small / large / context / tokens は 32 run の行で描き直した（記録の表の宣言は同じ key で追記してあり、全 run が揃った版で全行に戻す）。全 36 run は同じコミットで実行している。

\paragraph{反復上限と STE.}
図~\ref{fig:ste} と表~\ref{tab:small}、表~\ref{tab:large} に、上限ごとの平均 STE を示す。small では 3 モデルとも上限 1 から 40 まで STE が下がり続けた。DeepSeek は上限 1 の \airasval{deepseek-v4.1-flash-small-i01.trajectory_smape} から上限 20 で \airasval{deepseek-v4.1-flash-small-i20.trajectory_smape}、上限 40 で \airasval{deepseek-v4.1-flash-small-i40.trajectory_smape}、Qwen は \airasval{qwen3.8-27b-small-i01.trajectory_smape} から \airasval{qwen3.8-27b-small-i20.trajectory_smape}、\airasval{qwen3.8-27b-small-i40.trajectory_smape} で、どちらも上限 20 で論文 Table 1 の最良行 Gemini-2.5-Pro の \unverified{0.3212} \cite[s1.p5]{duan-2025-measuring} を下回った。Nemotron は上限 10 まで平坦（\airasval{nemotron-3-ultra-small-i01.trajectory_smape} から \airasval{nemotron-3-ultra-small-i10.trajectory_smape}）で、上限 20 で \airasval{nemotron-3-ultra-small-i20.trajectory_smape} に下がり、上限 40 は \airasval{nemotron-3-ultra-small-i40.trajectory_smape} と横ばい、上限 80 で \airasval{nemotron-3-ultra-small-i80.trajectory_smape} であった。

large では上限 20 までの改善が小さい。DeepSeek は上限 5 から 20 まで横ばい（\airasval{deepseek-v4.1-flash-large-i05.trajectory_smape}、\airasval{deepseek-v4.1-flash-large-i10.trajectory_smape}、\airasval{deepseek-v4.1-flash-large-i20.trajectory_smape}）だったが、上限 40 で \airasval{deepseek-v4.1-flash-large-i40.trajectory_smape} に下がった。Qwen も上限 20 の \airasval{qwen3.8-27b-large-i20.trajectory_smape} から上限 40 で \airasval{qwen3.8-27b-large-i40.trajectory_smape} に下がった。Nemotron は上限 20 の \airasval{nemotron-3-ultra-large-i20.trajectory_smape} から上限 40 で \airasval{nemotron-3-ultra-large-i40.trajectory_smape}、上限 80 で \airasval{nemotron-3-ultra-large-i80.trajectory_smape} と緩やかに下がった。

\begin{figure}[t]
  \centering
  \includegraphics[width=0.95\linewidth]{images/chart/ste_vs_budget_v2.png}
  \caption{反復上限ごとの平均 STE（低いほど良い）。左 small 137 件、右 large 213 件。点は取り込み済みの 32 run（DeepSeek と Qwen の上限 80 は未完）。横軸は対数。}
  \label{fig:ste}
\end{figure}

\begin{figure}[t]
  \centering
  \includegraphics[width=0.95\linewidth]{images/chart/iterations_vs_budget_v2.png}
  \caption{反復上限に対して実際に使った平均反復数（両軸対数）。対角線上なら予算を使い切り、下に外れれば上限前に自ら提出している。上限 20 までは 3 モデルとも使い切り、Nemotron は上限 40 以上で下に外れる。}
  \label{fig:it}
\end{figure}

{\scriptsize\setlength{\tabcolsep}{2.5pt}% AUTO-GENERATED by the update_record tool. DO NOT EDIT.
% verify_paper regenerates this file from record.json and the run outputs
% and fails on any difference, so manual edits always surface.
\begin{table}[t]
\centering
\caption{SciGym-small 137 件。行はモデル / 反復上限。平均 STE（低いほど良い）、Reaction Matching Score の F1（modifier なし・あり、高いほど良い）、有効な提出件数、提出までの平均反復数、上限前に自ら提出した件数、文脈長を超えて提出に至らなかった件数、提出に至らなかった件数（文脈超過・思考枯渇・3 試行とも打ち切りの合計）。 全 36 run（最終版）。}
\label{tab:small}
\begin{tabular}{lrrrrrrrr}
\hline
 & STE $\downarrow$ & F1 & F1$_m$ & n\_valid & 平均反復 & 早期提出 & 文脈超過 & 未提出 \\
\hline
DeepSeek-V4.1-Flash / 1 & \airasrecordlink[\#L442]{0.443} & \airasrecordlink[\#L429]{0.312} & \airasrecordlink[\#L431]{0.167} & \airasrecordlink[\#L445]{137} & \airasrecordlink[\#L452]{1.5} & \airasrecordlink[\#L453]{32} & \airasrecordlink[\#L454]{0} & \airasrecordlink[\#L456]{0} \\
DeepSeek-V4.1-Flash / 5 & \airasrecordlink[\#L529]{0.307} & \airasrecordlink[\#L516]{0.438} & \airasrecordlink[\#L518]{0.246} & \airasrecordlink[\#L532]{133} & \airasrecordlink[\#L539]{6.3} & \airasrecordlink[\#L540]{0} & \airasrecordlink[\#L541]{0} & \airasrecordlink[\#L543]{0} \\
DeepSeek-V4.1-Flash / 10 & \airasrecordlink[\#L616]{0.251} & \airasrecordlink[\#L603]{0.429} & \airasrecordlink[\#L605]{0.268} & \airasrecordlink[\#L619]{135} & \airasrecordlink[\#L626]{11.0} & \airasrecordlink[\#L627]{6} & \airasrecordlink[\#L628]{0} & \airasrecordlink[\#L630]{0} \\
DeepSeek-V4.1-Flash / 20 & \airasrecordlink[\#L703]{0.223} & \airasrecordlink[\#L690]{0.490} & \airasrecordlink[\#L692]{0.277} & \airasrecordlink[\#L706]{135} & \airasrecordlink[\#L713]{19.6} & \airasrecordlink[\#L714]{24} & \airasrecordlink[\#L715]{0} & \airasrecordlink[\#L717]{1} \\
DeepSeek-V4.1-Flash / 40 & \airasrecordlink[\#L7762]{0.198} & \airasrecordlink[\#L7749]{0.461} & \airasrecordlink[\#L7751]{0.295} & \airasrecordlink[\#L7765]{134} & \airasrecordlink[\#L7772]{33.2} & \airasrecordlink[\#L7773]{68} & \airasrecordlink[\#L7774]{0} & \airasrecordlink[\#L7776]{3} \\
DeepSeek-V4.1-Flash / 80 & \airasrecordlink[\#L790]{0.159} & \airasrecordlink[\#L777]{0.457} & \airasrecordlink[\#L779]{0.308} & \airasrecordlink[\#L793]{126} & \airasrecordlink[\#L800]{48.3} & \airasrecordlink[\#L801]{105} & \airasrecordlink[\#L802]{0} & \airasrecordlink[\#L804]{11} \\
Qwen3.8-27B / 1 & \airasrecordlink[\#L1370]{0.622} & \airasrecordlink[\#L1357]{0.108} & \airasrecordlink[\#L1359]{0.062} & \airasrecordlink[\#L1373]{133} & \airasrecordlink[\#L1380]{1.1} & \airasrecordlink[\#L1381]{77} & \airasrecordlink[\#L1382]{0} & \airasrecordlink[\#L1384]{0} \\
Qwen3.8-27B / 5 & \airasrecordlink[\#L1457]{0.372} & \airasrecordlink[\#L1444]{0.354} & \airasrecordlink[\#L1446]{0.192} & \airasrecordlink[\#L1460]{124} & \airasrecordlink[\#L1467]{6.6} & \airasrecordlink[\#L1468]{1} & \airasrecordlink[\#L1469]{0} & \airasrecordlink[\#L1471]{0} \\
Qwen3.8-27B / 10 & \airasrecordlink[\#L1544]{0.315} & \airasrecordlink[\#L1531]{0.366} & \airasrecordlink[\#L1533]{0.209} & \airasrecordlink[\#L1547]{125} & \airasrecordlink[\#L1554]{11.5} & \airasrecordlink[\#L1555]{4} & \airasrecordlink[\#L1556]{0} & \airasrecordlink[\#L1558]{1} \\
Qwen3.8-27B / 20 & \airasrecordlink[\#L1631]{0.263} & \airasrecordlink[\#L1618]{0.404} & \airasrecordlink[\#L1620]{0.239} & \airasrecordlink[\#L1634]{120} & \airasrecordlink[\#L1641]{19.9} & \airasrecordlink[\#L1642]{28} & \airasrecordlink[\#L1643]{0} & \airasrecordlink[\#L1645]{3} \\
Qwen3.8-27B / 40 & \airasrecordlink[\#L8008]{0.210} & \airasrecordlink[\#L7995]{0.482} & \airasrecordlink[\#L7997]{0.286} & \airasrecordlink[\#L8011]{125} & \airasrecordlink[\#L8018]{32.8} & \airasrecordlink[\#L8019]{69} & \airasrecordlink[\#L8020]{0} & \airasrecordlink[\#L8022]{7} \\
Qwen3.8-27B / 80 & \airasrecordlink[\#L1718]{0.166} & \airasrecordlink[\#L1705]{0.456} & \airasrecordlink[\#L1707]{0.268} & \airasrecordlink[\#L1721]{124} & \airasrecordlink[\#L1728]{48.0} & \airasrecordlink[\#L1729]{110} & \airasrecordlink[\#L1730]{0} & \airasrecordlink[\#L1732]{13} \\
Nemotron-3-Ultra / 1 & \airasrecordlink[\#L2298]{0.392} & \airasrecordlink[\#L2285]{0.369} & \airasrecordlink[\#L2287]{0.201} & \airasrecordlink[\#L2301]{135} & \airasrecordlink[\#L2308]{2.4} & \airasrecordlink[\#L2309]{0} & \airasrecordlink[\#L2310]{0} & \airasrecordlink[\#L2312]{0} \\
Nemotron-3-Ultra / 5 & \airasrecordlink[\#L2385]{0.375} & \airasrecordlink[\#L2372]{0.424} & \airasrecordlink[\#L2374]{0.243} & \airasrecordlink[\#L2388]{133} & \airasrecordlink[\#L2395]{6.4} & \airasrecordlink[\#L2396]{1} & \airasrecordlink[\#L2397]{0} & \airasrecordlink[\#L2399]{0} \\
Nemotron-3-Ultra / 10 & \airasrecordlink[\#L2472]{0.369} & \airasrecordlink[\#L2459]{0.421} & \airasrecordlink[\#L2461]{0.241} & \airasrecordlink[\#L2475]{132} & \airasrecordlink[\#L2482]{10.9} & \airasrecordlink[\#L2483]{10} & \airasrecordlink[\#L2484]{0} & \airasrecordlink[\#L2486]{0} \\
Nemotron-3-Ultra / 20 & \airasrecordlink[\#L2559]{0.283} & \airasrecordlink[\#L2546]{0.457} & \airasrecordlink[\#L2548]{0.262} & \airasrecordlink[\#L2562]{137} & \airasrecordlink[\#L2569]{18.5} & \airasrecordlink[\#L2570]{49} & \airasrecordlink[\#L2571]{0} & \airasrecordlink[\#L2573]{0} \\
Nemotron-3-Ultra / 40 & \airasrecordlink[\#L8254]{0.282} & \airasrecordlink[\#L8241]{0.401} & \airasrecordlink[\#L8243]{0.242} & \airasrecordlink[\#L8257]{131} & \airasrecordlink[\#L8264]{26.4} & \airasrecordlink[\#L8265]{114} & \airasrecordlink[\#L8266]{0} & \airasrecordlink[\#L8268]{0} \\
Nemotron-3-Ultra / 80 & \airasrecordlink[\#L2646]{0.255} & \airasrecordlink[\#L2633]{0.404} & \airasrecordlink[\#L2635]{0.229} & \airasrecordlink[\#L2649]{137} & \airasrecordlink[\#L2656]{31.9} & \airasrecordlink[\#L2657]{134} & \airasrecordlink[\#L2658]{0} & \airasrecordlink[\#L2660]{0} \\
\hline
\end{tabular}
\end{table}
}

{\scriptsize\setlength{\tabcolsep}{2.5pt}% AUTO-GENERATED by the update_record tool. DO NOT EDIT.
% verify_paper regenerates this file from record.json and the run outputs
% and fails on any difference, so manual edits always surface.
\begin{table}[t]
\centering
\caption{SciGym-large 213 件。列は表 small と同じ。 全 36 run（最終版）。}
\label{tab:large}
\begin{tabular}{lrrrrrrrr}
\hline
 & STE $\downarrow$ & F1 & F1$_m$ & n\_valid & 平均反復 & 早期提出 & 文脈超過 & 未提出 \\
\hline
DeepSeek-V4.1-Flash / 1 & \airasrecordlink[\#L902]{0.508} & \airasrecordlink[\#L890]{0.131} & \airasrecordlink[\#L891]{0.063} & \airasrecordlink[\#L904]{209} & \airasrecordlink[\#L911]{1.6} & \airasrecordlink[\#L912]{40} & \airasrecordlink[\#L913]{0} & \airasrecordlink[\#L915]{0} \\
DeepSeek-V4.1-Flash / 5 & \airasrecordlink[\#L985]{0.438} & \airasrecordlink[\#L973]{0.229} & \airasrecordlink[\#L974]{0.108} & \airasrecordlink[\#L987]{211} & \airasrecordlink[\#L994]{6.2} & \airasrecordlink[\#L995]{0} & \airasrecordlink[\#L996]{0} & \airasrecordlink[\#L998]{0} \\
DeepSeek-V4.1-Flash / 10 & \airasrecordlink[\#L1068]{0.437} & \airasrecordlink[\#L1056]{0.212} & \airasrecordlink[\#L1057]{0.105} & \airasrecordlink[\#L1070]{208} & \airasrecordlink[\#L1077]{11.1} & \airasrecordlink[\#L1078]{0} & \airasrecordlink[\#L1079]{1} & \airasrecordlink[\#L1081]{1} \\
DeepSeek-V4.1-Flash / 20 & \airasrecordlink[\#L1151]{0.443} & \airasrecordlink[\#L1139]{0.212} & \airasrecordlink[\#L1140]{0.117} & \airasrecordlink[\#L1153]{208} & \airasrecordlink[\#L1160]{20.9} & \airasrecordlink[\#L1161]{7} & \airasrecordlink[\#L1162]{3} & \airasrecordlink[\#L1164]{3} \\
DeepSeek-V4.1-Flash / 40 & \airasrecordlink[\#L7873]{0.352} & \airasrecordlink[\#L7861]{0.249} & \airasrecordlink[\#L7862]{0.135} & \airasrecordlink[\#L7875]{208} & \airasrecordlink[\#L7882]{39.2} & \airasrecordlink[\#L7883]{56} & \airasrecordlink[\#L7884]{3} & \airasrecordlink[\#L7886]{5} \\
DeepSeek-V4.1-Flash / 80 & \airasrecordlink[\#L1234]{0.293} & \airasrecordlink[\#L1222]{0.258} & \airasrecordlink[\#L1223]{0.142} & \airasrecordlink[\#L1236]{205} & \airasrecordlink[\#L1243]{65.4} & \airasrecordlink[\#L1244]{158} & \airasrecordlink[\#L1245]{3} & \airasrecordlink[\#L1247]{8} \\
Qwen3.8-27B / 1 & \airasrecordlink[\#L1830]{0.582} & \airasrecordlink[\#L1818]{0.072} & \airasrecordlink[\#L1819]{0.033} & \airasrecordlink[\#L1832]{200} & \airasrecordlink[\#L1839]{1.4} & \airasrecordlink[\#L1840]{95} & \airasrecordlink[\#L1841]{0} & \airasrecordlink[\#L1843]{0} \\
Qwen3.8-27B / 5 & \airasrecordlink[\#L1913]{0.483} & \airasrecordlink[\#L1901]{0.161} & \airasrecordlink[\#L1902]{0.079} & \airasrecordlink[\#L1915]{188} & \airasrecordlink[\#L1922]{6.9} & \airasrecordlink[\#L1923]{0} & \airasrecordlink[\#L1924]{0} & \airasrecordlink[\#L1926]{0} \\
Qwen3.8-27B / 10 & \airasrecordlink[\#L1996]{0.471} & \airasrecordlink[\#L1984]{0.164} & \airasrecordlink[\#L1985]{0.089} & \airasrecordlink[\#L1998]{195} & \airasrecordlink[\#L2005]{11.7} & \airasrecordlink[\#L2006]{0} & \airasrecordlink[\#L2007]{0} & \airasrecordlink[\#L2009]{0} \\
Qwen3.8-27B / 20 & \airasrecordlink[\#L2079]{0.431} & \airasrecordlink[\#L2067]{0.210} & \airasrecordlink[\#L2068]{0.112} & \airasrecordlink[\#L2081]{193} & \airasrecordlink[\#L2088]{21.2} & \airasrecordlink[\#L2089]{9} & \airasrecordlink[\#L2090]{1} & \airasrecordlink[\#L2092]{1} \\
Qwen3.8-27B / 40 & \airasrecordlink[\#L8119]{0.373} & \airasrecordlink[\#L8107]{0.244} & \airasrecordlink[\#L8108]{0.142} & \airasrecordlink[\#L8121]{200} & \airasrecordlink[\#L8128]{39.5} & \airasrecordlink[\#L8129]{51} & \airasrecordlink[\#L8130]{4} & \airasrecordlink[\#L8132]{8} \\
Qwen3.8-27B / 80 & \airasrecordlink[\#L2162]{0.300} & \airasrecordlink[\#L2150]{0.287} & \airasrecordlink[\#L2151]{0.163} & \airasrecordlink[\#L2164]{195} & \airasrecordlink[\#L2171]{67.2} & \airasrecordlink[\#L2172]{143} & \airasrecordlink[\#L2173]{5} & \airasrecordlink[\#L2175]{15} \\
Nemotron-3-Ultra / 1 & \airasrecordlink[\#L2758]{0.516} & \airasrecordlink[\#L2746]{0.168} & \airasrecordlink[\#L2747]{0.083} & \airasrecordlink[\#L2760]{207} & \airasrecordlink[\#L2767]{2.6} & \airasrecordlink[\#L2768]{0} & \airasrecordlink[\#L2769]{0} & \airasrecordlink[\#L2771]{0} \\
Nemotron-3-Ultra / 5 & \airasrecordlink[\#L2841]{0.521} & \airasrecordlink[\#L2829]{0.193} & \airasrecordlink[\#L2830]{0.103} & \airasrecordlink[\#L2843]{203} & \airasrecordlink[\#L2850]{6.5} & \airasrecordlink[\#L2851]{0} & \airasrecordlink[\#L2852]{1} & \airasrecordlink[\#L2854]{1} \\
Nemotron-3-Ultra / 10 & \airasrecordlink[\#L2924]{0.488} & \airasrecordlink[\#L2912]{0.218} & \airasrecordlink[\#L2913]{0.123} & \airasrecordlink[\#L2926]{204} & \airasrecordlink[\#L2933]{11.1} & \airasrecordlink[\#L2934]{8} & \airasrecordlink[\#L2935]{3} & \airasrecordlink[\#L2937]{3} \\
Nemotron-3-Ultra / 20 & \airasrecordlink[\#L3007]{0.418} & \airasrecordlink[\#L2995]{0.218} & \airasrecordlink[\#L2996]{0.112} & \airasrecordlink[\#L3009]{200} & \airasrecordlink[\#L3016]{20.0} & \airasrecordlink[\#L3017]{60} & \airasrecordlink[\#L3018]{3} & \airasrecordlink[\#L3020]{3} \\
Nemotron-3-Ultra / 40 & \airasrecordlink[\#L8365]{0.392} & \airasrecordlink[\#L8353]{0.226} & \airasrecordlink[\#L8354]{0.128} & \airasrecordlink[\#L8367]{205} & \airasrecordlink[\#L8374]{30.0} & \airasrecordlink[\#L8375]{170} & \airasrecordlink[\#L8376]{3} & \airasrecordlink[\#L8378]{3} \\
Nemotron-3-Ultra / 80 & \airasrecordlink[\#L3090]{0.377} & \airasrecordlink[\#L3078]{0.213} & \airasrecordlink[\#L3079]{0.110} & \airasrecordlink[\#L3092]{207} & \airasrecordlink[\#L3099]{38.8} & \airasrecordlink[\#L3100]{194} & \airasrecordlink[\#L3101]{5} & \airasrecordlink[\#L3103]{5} \\
\hline
\end{tabular}
\end{table}
}

\paragraph{主張の判定.}
Nemotron-3-Ultra の 4 主張が判定された。Claim 5（small、上限 80 が上限 20 より 0.03 以上良い）は、差が \airasval{nemotron-3-ultra-small-i80.trajectory_smape} $-$ \airasval{nemotron-3-ultra-small-i20.trajectory_smape} $= -0.028$ で margin $-0.03$ にわずかに届かず\textbf{反証}された。Claim 6（large、同じ基準）は差 $-0.041$ で\textbf{支持}された。Claim 11・12（上限 40 から 80 への改善が 0.05 未満）は、small で差 $0.026$、large で差 $0.015$ と、どちらも\textbf{支持}された。差は予測区間（Claim 5: $[-0.12, -0.02]$、Claim 6: 同、Claim 11: $[-0.01, 0.06]$、Claim 12: 同）の内側である。DeepSeek と Qwen の 8 主張は上限 80 の run の完了後に判定する。

\paragraph{予算を使うか.}
図~\ref{fig:it} に、上限に対して実際に使った平均反復数を示す。上限 20 までは 3 モデルとも平均反復数が上限とほぼ一致し、予算を使い切ってから提出している。上限 40 では分かれる。DeepSeek と Qwen は平均 \airasval{deepseek-v4.1-flash-small-i40.mean_iterations}（small）、\airasval{deepseek-v4.1-flash-large-i40.mean_iterations}（large）、Qwen は \airasval{qwen3.8-27b-small-i40.mean_iterations}、\airasval{qwen3.8-27b-large-i40.mean_iterations} と大半を使うのに対し、Nemotron は small で \airasval{nemotron-3-ultra-small-i40.mean_iterations}、large で \airasval{nemotron-3-ultra-large-i40.mean_iterations} にとどまり、上限 80 でも small \airasval{nemotron-3-ultra-small-i80.mean_iterations}、large \airasval{nemotron-3-ultra-large-i80.mean_iterations} である。上限前に自ら提出した件数は、Nemotron の small 上限 80 で 137 件中 \airasval{nemotron-3-ultra-small-i80.n_early_submit} 件、large 上限 80 で 213 件中 \airasval{nemotron-3-ultra-large-i80.n_early_submit} 件に達する。上限 1 では、Qwen が small の \airasval{qwen3.8-27b-small-i01.n_early_submit} 件、large の \airasval{qwen3.8-27b-large-i01.n_early_submit} 件で実験を行わずに提出した。平均反復数は無効な提出のあとのデバッグ反復を含むため、上限を超えることがある。

\paragraph{有効な提出と文脈.}
有効な SBML が提出された件数（表の n\_valid）は、DeepSeek と Nemotron が small で 131〜137 件、Qwen は上限 40 で \airasval{qwen3.8-27b-small-i40.n_valid_submissions} 件と最も少ない。表~\ref{tab:context} に文脈の使用量を示す。文脈長を 1M に揃えたことで、文脈超過で提出に至らなかった件は 32 run 合計 \unverified{30} 件（すべて large。上限 40 で DeepSeek \airasval{deepseek-v4.1-flash-large-i40.n_context_overflow} 件、Qwen \airasval{qwen3.8-27b-large-i40.n_context_overflow} 件、Nemotron の上限 80 で \airasval{nemotron-3-ultra-large-i80.n_context_overflow} 件）にとどまった。1 呼び出しあたりの最大プロンプト長は large の上限 10 以上で 100 万トークン近くに達している。思考だけで出力上限を使い切って本文が出なかった件（n\_thinking\_exhausted）は全 32 run で 0 であった。3 試行とも打ち切られて提出に至らなかった件は、Qwen の上限 40 で small \unverified{7} 件、large \unverified{8} 件、DeepSeek の上限 40 で small \unverified{3} 件、large \unverified{5} 件あり、表の n\_valid の差の一部を占める。

{\tiny\setlength{\tabcolsep}{2pt}% AUTO-GENERATED by the update_record tool. DO NOT EDIT.
% verify_paper regenerates this file from record.json and the run outputs
% and fails on any difference, so manual edits always surface.
\begin{table}[t]
\centering
\caption{文脈の使用量。行はモデル split / 反復上限（DeepSeek = DeepSeek-V4.1-Flash、Qwen = Qwen3.8-27B、Nemotron = Nemotron-3-Ultra）。反復上限、サーバーの文脈長設定、LLM 呼び出し数、1 呼び出しあたりの平均プロンプト長と最大プロンプト長（トークン）、文脈超過で提出に至らなかった件数。 実行中の 4 run（DeepSeek と Qwen の上限 80）は含めない（暫定版 v2）。}
\label{tab:context}
\begin{tabular}{lrrrrrr}
\hline
 & 上限 & 文脈長 & 呼び出し & 平均プロンプト & 最大プロンプト & 文脈超過 \\
\hline
DeepSeek small / 1 & \airasrecordlink[\#L435]{1} & \airasrecordlink[\#L436]{1048576} & \airasrecordlink[\#L426]{346} & \airasrecordlink[\#L429]{3683} & \airasrecordlink[\#L428]{14526} & \airasrecordlink[\#L432]{0} \\
DeepSeek small / 5 & \airasrecordlink[\#L522]{5} & \airasrecordlink[\#L523]{1048576} & \airasrecordlink[\#L513]{998} & \airasrecordlink[\#L516]{6905} & \airasrecordlink[\#L515]{26663} & \airasrecordlink[\#L519]{0} \\
DeepSeek small / 10 & \airasrecordlink[\#L609]{10} & \airasrecordlink[\#L610]{1048576} & \airasrecordlink[\#L600]{1643} & \airasrecordlink[\#L603]{10243} & \airasrecordlink[\#L602]{35509} & \airasrecordlink[\#L606]{0} \\
DeepSeek small / 20 & \airasrecordlink[\#L696]{20} & \airasrecordlink[\#L697]{1048576} & \airasrecordlink[\#L687]{2826} & \airasrecordlink[\#L690]{16390} & \airasrecordlink[\#L689]{63174} & \airasrecordlink[\#L693]{0} \\
DeepSeek small / 40 & \airasrecordlink[\#L6240]{40} & \airasrecordlink[\#L6241]{1048576} & \airasrecordlink[\#L6231]{4660} & \airasrecordlink[\#L6234]{28311} & \airasrecordlink[\#L6233]{122867} & \airasrecordlink[\#L6237]{0} \\
Qwen small / 1 & \airasrecordlink[\#L1203]{1} & \airasrecordlink[\#L1204]{1000000} & \airasrecordlink[\#L1194]{290} & \airasrecordlink[\#L1197]{4213} & \airasrecordlink[\#L1196]{62449} & \airasrecordlink[\#L1200]{0} \\
Qwen small / 5 & \airasrecordlink[\#L1290]{5} & \airasrecordlink[\#L1291]{1000000} & \airasrecordlink[\#L1281]{1041} & \airasrecordlink[\#L1284]{10128} & \airasrecordlink[\#L1283]{78199} & \airasrecordlink[\#L1287]{0} \\
Qwen small / 10 & \airasrecordlink[\#L1377]{10} & \airasrecordlink[\#L1378]{1000000} & \airasrecordlink[\#L1368]{1706} & \airasrecordlink[\#L1371]{14767} & \airasrecordlink[\#L1370]{53460} & \airasrecordlink[\#L1374]{0} \\
Qwen small / 20 & \airasrecordlink[\#L1464]{20} & \airasrecordlink[\#L1465]{1000000} & \airasrecordlink[\#L1455]{2827} & \airasrecordlink[\#L1458]{22927} & \airasrecordlink[\#L1457]{88692} & \airasrecordlink[\#L1461]{0} \\
Qwen small / 40 & \airasrecordlink[\#L6482]{40} & \airasrecordlink[\#L6483]{1000000} & \airasrecordlink[\#L6473]{4506} & \airasrecordlink[\#L6476]{41868} & \airasrecordlink[\#L6475]{322306} & \airasrecordlink[\#L6479]{0} \\
Nemotron small / 1 & \airasrecordlink[\#L1971]{1} & \airasrecordlink[\#L1972]{1048576} & \airasrecordlink[\#L1962]{462} & \airasrecordlink[\#L1965]{5327} & \airasrecordlink[\#L1964]{26211} & \airasrecordlink[\#L1968]{0} \\
Nemotron small / 5 & \airasrecordlink[\#L2058]{5} & \airasrecordlink[\#L2059]{1048576} & \airasrecordlink[\#L2049]{1010} & \airasrecordlink[\#L2052]{10560} & \airasrecordlink[\#L2051]{44190} & \airasrecordlink[\#L2055]{0} \\
Nemotron small / 10 & \airasrecordlink[\#L2145]{10} & \airasrecordlink[\#L2146]{1048576} & \airasrecordlink[\#L2136]{1631} & \airasrecordlink[\#L2139]{18264} & \airasrecordlink[\#L2138]{79051} & \airasrecordlink[\#L2142]{0} \\
Nemotron small / 20 & \airasrecordlink[\#L2232]{20} & \airasrecordlink[\#L2233]{1048576} & \airasrecordlink[\#L2223]{2665} & \airasrecordlink[\#L2226]{32090} & \airasrecordlink[\#L2225]{142060} & \airasrecordlink[\#L2229]{0} \\
Nemotron small / 40 & \airasrecordlink[\#L6724]{40} & \airasrecordlink[\#L6725]{1048576} & \airasrecordlink[\#L6715]{3759} & \airasrecordlink[\#L6718]{49467} & \airasrecordlink[\#L6717]{245901} & \airasrecordlink[\#L6721]{0} \\
Nemotron small / 80 & \airasrecordlink[\#L2319]{80} & \airasrecordlink[\#L2320]{1048576} & \airasrecordlink[\#L2310]{4504} & \airasrecordlink[\#L2313]{72465} & \airasrecordlink[\#L2312]{485039} & \airasrecordlink[\#L2316]{0} \\
DeepSeek large / 1 & \airasrecordlink[\#L823]{1} & \airasrecordlink[\#L824]{1048576} & \airasrecordlink[\#L814]{554} & \airasrecordlink[\#L817]{9259} & \airasrecordlink[\#L816]{187265} & \airasrecordlink[\#L820]{0} \\
DeepSeek large / 5 & \airasrecordlink[\#L906]{5} & \airasrecordlink[\#L907]{1048576} & \airasrecordlink[\#L897]{1536} & \airasrecordlink[\#L900]{19247} & \airasrecordlink[\#L899]{769628} & \airasrecordlink[\#L903]{0} \\
DeepSeek large / 10 & \airasrecordlink[\#L989]{10} & \airasrecordlink[\#L990]{1048576} & \airasrecordlink[\#L980]{2568} & \airasrecordlink[\#L983]{26376} & \airasrecordlink[\#L982]{993655} & \airasrecordlink[\#L986]{1} \\
DeepSeek large / 20 & \airasrecordlink[\#L1072]{20} & \airasrecordlink[\#L1073]{1048576} & \airasrecordlink[\#L1063]{4629} & \airasrecordlink[\#L1066]{42082} & \airasrecordlink[\#L1065]{1003501} & \airasrecordlink[\#L1069]{3} \\
DeepSeek large / 40 & \airasrecordlink[\#L6359]{40} & \airasrecordlink[\#L6360]{1048576} & \airasrecordlink[\#L6350]{8478} & \airasrecordlink[\#L6353]{61628} & \airasrecordlink[\#L6352]{1006914} & \airasrecordlink[\#L6356]{3} \\
Qwen large / 1 & \airasrecordlink[\#L1591]{1} & \airasrecordlink[\#L1592]{1000000} & \airasrecordlink[\#L1582]{505} & \airasrecordlink[\#L1585]{10596} & \airasrecordlink[\#L1584]{226549} & \airasrecordlink[\#L1588]{0} \\
Qwen large / 5 & \airasrecordlink[\#L1674]{5} & \airasrecordlink[\#L1675]{1000000} & \airasrecordlink[\#L1665]{1674} & \airasrecordlink[\#L1668]{26205} & \airasrecordlink[\#L1667]{632543} & \airasrecordlink[\#L1671]{0} \\
Qwen large / 10 & \airasrecordlink[\#L1757]{10} & \airasrecordlink[\#L1758]{1000000} & \airasrecordlink[\#L1748]{2701} & \airasrecordlink[\#L1751]{34096} & \airasrecordlink[\#L1750]{778836} & \airasrecordlink[\#L1754]{0} \\
Qwen large / 20 & \airasrecordlink[\#L1840]{20} & \airasrecordlink[\#L1841]{1000000} & \airasrecordlink[\#L1831]{4722} & \airasrecordlink[\#L1834]{49151} & \airasrecordlink[\#L1833]{863229} & \airasrecordlink[\#L1837]{1} \\
Qwen large / 40 & \airasrecordlink[\#L6601]{40} & \airasrecordlink[\#L6602]{1000000} & \airasrecordlink[\#L6592]{8431} & \airasrecordlink[\#L6595]{76681} & \airasrecordlink[\#L6594]{933789} & \airasrecordlink[\#L6598]{4} \\
Nemotron large / 1 & \airasrecordlink[\#L2441]{1} & \airasrecordlink[\#L2442]{1048576} & \airasrecordlink[\#L2432]{764} & \airasrecordlink[\#L2435]{16584} & \airasrecordlink[\#L2434]{505482} & \airasrecordlink[\#L2438]{0} \\
Nemotron large / 5 & \airasrecordlink[\#L2524]{5} & \airasrecordlink[\#L2525]{1048576} & \airasrecordlink[\#L2515]{1598} & \airasrecordlink[\#L2518]{33418} & \airasrecordlink[\#L2517]{983815} & \airasrecordlink[\#L2521]{1} \\
Nemotron large / 10 & \airasrecordlink[\#L2607]{10} & \airasrecordlink[\#L2608]{1048576} & \airasrecordlink[\#L2598]{2559} & \airasrecordlink[\#L2601]{46983} & \airasrecordlink[\#L2600]{1005002} & \airasrecordlink[\#L2604]{3} \\
Nemotron large / 20 & \airasrecordlink[\#L2690]{20} & \airasrecordlink[\#L2691]{1048576} & \airasrecordlink[\#L2681]{4437} & \airasrecordlink[\#L2684]{69866} & \airasrecordlink[\#L2683]{985351} & \airasrecordlink[\#L2687]{3} \\
Nemotron large / 40 & \airasrecordlink[\#L6845]{40} & \airasrecordlink[\#L6846]{1048576} & \airasrecordlink[\#L6836]{6520} & \airasrecordlink[\#L6839]{95004} & \airasrecordlink[\#L6838]{988916} & \airasrecordlink[\#L6842]{3} \\
Nemotron large / 80 & \airasrecordlink[\#L2773]{80} & \airasrecordlink[\#L2774]{1048576} & \airasrecordlink[\#L2764]{8405} & \airasrecordlink[\#L2767]{140176} & \airasrecordlink[\#L2766]{1013986} & \airasrecordlink[\#L2770]{5} \\
\hline
\end{tabular}
\end{table}
}

\paragraph{トークン.}
表~\ref{tab:tokens} に入力・出力トークン総数を示す。入力トークンは反復とともに会話履歴が積み上がるため上限に対して超線形に増える。出力トークンには思考を含み、Qwen は他の 2 モデルより 2〜3 倍多い。思考が本文に混じって除去した応答は全 run で 0 件である。

{\scriptsize\setlength{\tabcolsep}{2.5pt}% AUTO-GENERATED by the update_record tool. DO NOT EDIT.
% verify_paper regenerates this file from record.json and the run outputs
% and fails on any difference, so manual edits always surface.
\begin{table}[t]
\centering
\caption{トークン使用量。行はモデル split / 反復上限（DeepSeek = DeepSeek-V4.1-Flash、Qwen = Qwen3.8-27B、Nemotron = Nemotron-3-Ultra）。入力・出力トークン総数、LLM 呼び出し数、思考が本文に混じって除去した応答数。 全 36 run（最終版）。}
\label{tab:tokens}
\begin{tabular}{lrrrr}
\hline
 & 入力 & 出力 & 呼び出し & 思考混入 \\
\hline
DeepSeek small / 1 & \airasrecordlink[\#L435]{1274290} & \airasrecordlink[\#L436]{2297420} & \airasrecordlink[\#L437]{346} & \airasrecordlink[\#L438]{3} \\
DeepSeek small / 5 & \airasrecordlink[\#L522]{6891226} & \airasrecordlink[\#L523]{3364660} & \airasrecordlink[\#L524]{998} & \airasrecordlink[\#L525]{2} \\
DeepSeek small / 10 & \airasrecordlink[\#L609]{16828516} & \airasrecordlink[\#L610]{4474220} & \airasrecordlink[\#L611]{1643} & \airasrecordlink[\#L612]{4} \\
DeepSeek small / 20 & \airasrecordlink[\#L696]{46317718} & \airasrecordlink[\#L697]{6966753} & \airasrecordlink[\#L698]{2826} & \airasrecordlink[\#L699]{7} \\
DeepSeek small / 40 & \airasrecordlink[\#L7755]{131929994} & \airasrecordlink[\#L7756]{12322098} & \airasrecordlink[\#L7757]{4660} & \airasrecordlink[\#L7758]{79} \\
DeepSeek small / 80 & \airasrecordlink[\#L783]{273180219} & \airasrecordlink[\#L784]{16916137} & \airasrecordlink[\#L785]{6579} & \airasrecordlink[\#L786]{137} \\
Qwen small / 1 & \airasrecordlink[\#L1363]{1221631} & \airasrecordlink[\#L1364]{4386426} & \airasrecordlink[\#L1365]{290} & \airasrecordlink[\#L1366]{285} \\
Qwen small / 5 & \airasrecordlink[\#L1450]{10542963} & \airasrecordlink[\#L1451]{11869691} & \airasrecordlink[\#L1452]{1041} & \airasrecordlink[\#L1453]{1037} \\
Qwen small / 10 & \airasrecordlink[\#L1537]{25192969} & \airasrecordlink[\#L1538]{14662945} & \airasrecordlink[\#L1539]{1706} & \airasrecordlink[\#L1540]{1703} \\
Qwen small / 20 & \airasrecordlink[\#L1624]{64815651} & \airasrecordlink[\#L1625]{22112889} & \airasrecordlink[\#L1626]{2827} & \airasrecordlink[\#L1627]{2815} \\
Qwen small / 40 & \airasrecordlink[\#L8001]{188658839} & \airasrecordlink[\#L8002]{31328832} & \airasrecordlink[\#L8003]{4506} & \airasrecordlink[\#L8004]{4421} \\
Qwen small / 80 & \airasrecordlink[\#L1711]{477679860} & \airasrecordlink[\#L1712]{42977847} & \airasrecordlink[\#L1713]{6563} & \airasrecordlink[\#L1714]{6260} \\
Nemotron small / 1 & \airasrecordlink[\#L2291]{2460942} & \airasrecordlink[\#L2292]{1454084} & \airasrecordlink[\#L2293]{462} & \airasrecordlink[\#L2294]{18} \\
Nemotron small / 5 & \airasrecordlink[\#L2378]{10665914} & \airasrecordlink[\#L2379]{2619710} & \airasrecordlink[\#L2380]{1010} & \airasrecordlink[\#L2381]{68} \\
Nemotron small / 10 & \airasrecordlink[\#L2465]{29788521} & \airasrecordlink[\#L2466]{4196975} & \airasrecordlink[\#L2467]{1631} & \airasrecordlink[\#L2468]{104} \\
Nemotron small / 20 & \airasrecordlink[\#L2552]{85518859} & \airasrecordlink[\#L2553]{7454118} & \airasrecordlink[\#L2554]{2665} & \airasrecordlink[\#L2555]{195} \\
Nemotron small / 40 & \airasrecordlink[\#L8247]{185946822} & \airasrecordlink[\#L8248]{10258370} & \airasrecordlink[\#L8249]{3759} & \airasrecordlink[\#L8250]{300} \\
Nemotron small / 80 & \airasrecordlink[\#L2639]{326380441} & \airasrecordlink[\#L2640]{13559568} & \airasrecordlink[\#L2641]{4504} & \airasrecordlink[\#L2642]{248} \\
DeepSeek large / 1 & \airasrecordlink[\#L905]{5129492} & \airasrecordlink[\#L906]{3829635} & \airasrecordlink[\#L907]{554} & \airasrecordlink[\#L908]{1} \\
DeepSeek large / 5 & \airasrecordlink[\#L988]{29564117} & \airasrecordlink[\#L989]{4944205} & \airasrecordlink[\#L990]{1536} & \airasrecordlink[\#L991]{4} \\
DeepSeek large / 10 & \airasrecordlink[\#L1071]{67733790} & \airasrecordlink[\#L1072]{6907247} & \airasrecordlink[\#L1073]{2568} & \airasrecordlink[\#L1074]{10} \\
DeepSeek large / 20 & \airasrecordlink[\#L1154]{194797668} & \airasrecordlink[\#L1155]{11480927} & \airasrecordlink[\#L1156]{4629} & \airasrecordlink[\#L1157]{30} \\
DeepSeek large / 40 & \airasrecordlink[\#L7876]{522485256} & \airasrecordlink[\#L7877]{20496954} & \airasrecordlink[\#L7878]{8478} & \airasrecordlink[\#L7879]{121} \\
DeepSeek large / 80 & \airasrecordlink[\#L1237]{1223909791} & \airasrecordlink[\#L1238]{32242354} & \airasrecordlink[\#L1239]{13860} & \airasrecordlink[\#L1240]{457} \\
Qwen large / 1 & \airasrecordlink[\#L1833]{5350897} & \airasrecordlink[\#L1834]{8630323} & \airasrecordlink[\#L1835]{505} & \airasrecordlink[\#L1836]{482} \\
Qwen large / 5 & \airasrecordlink[\#L1916]{43866762} & \airasrecordlink[\#L1917]{18209921} & \airasrecordlink[\#L1918]{1674} & \airasrecordlink[\#L1919]{1642} \\
Qwen large / 10 & \airasrecordlink[\#L1999]{92094533} & \airasrecordlink[\#L2000]{23167102} & \airasrecordlink[\#L2001]{2701} & \airasrecordlink[\#L2002]{2654} \\
Qwen large / 20 & \airasrecordlink[\#L2082]{232090863} & \airasrecordlink[\#L2083]{35409870} & \airasrecordlink[\#L2084]{4722} & \airasrecordlink[\#L2085]{4630} \\
Qwen large / 40 & \airasrecordlink[\#L8122]{646495199} & \airasrecordlink[\#L8123]{55672070} & \airasrecordlink[\#L8124]{8431} & \airasrecordlink[\#L8125]{8091} \\
Qwen large / 80 & \airasrecordlink[\#L2165]{1844165376} & \airasrecordlink[\#L2166]{82746596} & \airasrecordlink[\#L2167]{14017} & \airasrecordlink[\#L2168]{12969} \\
Nemotron large / 1 & \airasrecordlink[\#L2761]{12669803} & \airasrecordlink[\#L2762]{2800896} & \airasrecordlink[\#L2763]{764} & \airasrecordlink[\#L2764]{24} \\
Nemotron large / 5 & \airasrecordlink[\#L2844]{53401416} & \airasrecordlink[\#L2845]{4593209} & \airasrecordlink[\#L2846]{1598} & \airasrecordlink[\#L2847]{107} \\
Nemotron large / 10 & \airasrecordlink[\#L2927]{120229323} & \airasrecordlink[\#L2928]{6704393} & \airasrecordlink[\#L2929]{2559} & \airasrecordlink[\#L2930]{122} \\
Nemotron large / 20 & \airasrecordlink[\#L3010]{309996279} & \airasrecordlink[\#L3011]{12988712} & \airasrecordlink[\#L3012]{4437} & \airasrecordlink[\#L3013]{278} \\
Nemotron large / 40 & \airasrecordlink[\#L8368]{619428562} & \airasrecordlink[\#L8369]{19066662} & \airasrecordlink[\#L8370]{6520} & \airasrecordlink[\#L8371]{510} \\
Nemotron large / 80 & \airasrecordlink[\#L3093]{1178177246} & \airasrecordlink[\#L3094]{25610557} & \airasrecordlink[\#L3095]{8405} & \airasrecordlink[\#L3096]{512} \\
\hline
\end{tabular}
\end{table}
}

\subsection{補助分析: 予算の使い方（記録外）}
\label{sec:behavior}
以下は各件の会話履歴（\texttt{chat\_history.yaml}）と実験要求の記録（\texttt{experiment\_request.json}）から数えた値で、記録の検証対象ではない（すべて \texttt{unverified}）。行動の分類は Controller と同じ見出し（\texttt{\#\#\# Experiment}、\texttt{\#\#\# Code}、\texttt{final\_sbml} の代入）で行った。件ごとの STE は Controller が公式 Evaluator で採点した \texttt{evaluation.json} の値を用いた。

\paragraph{増えた予算はコードに使われる.}
表~\ref{tab:behavior} に、1 件あたりの実験要求回数とコード実行回数、提出した反復の中央値、摂動を伴わない実験要求の割合を示す。上限を 20 から 40 に倍にしても、実験要求は DeepSeek で \unverified{5.0} から \unverified{7.8} 回（small）、Qwen で \unverified{2.3} から \unverified{2.7} 回、Nemotron で \unverified{4.7} から \unverified{5.2} 回とわずかに増えるだけで、増分の大半はコード実行（small の上限 40 で DeepSeek \unverified{24.6} 回、Qwen \unverified{25.6} 回、Nemotron \unverified{17.8} 回）に使われている。Qwen は実験が最も少なく、コード中心の戦略である。軌道上の位置で見ると、実験要求は前 1/3 に集中し（上限 20 の small で、前 1/3 の行動のうち実験が DeepSeek \unverified{40}\%、Qwen \unverified{26}\%、Nemotron \unverified{47}\%）、後 1/3 ではコードがほとんどになる（同 \unverified{12}\%、\unverified{5}\%、\unverified{8}\%）。摂動を伴わない実験要求（初期条件のまま時系列を観測する）は上限 1 では 9 割以上、上限 20 以上でも 2〜4 割を占め、最初の実験は「素の時系列を見る」ことに使われている。提出した反復の中央値は、Nemotron が上限 40 でも 80 でも small で \unverified{21}、large で \unverified{25}〜\unverified{27} であり、上限に関わらず 20 前後で切り上げている。

\begin{table}[t]
\centering
\caption{予算の使い方（記録外の補助分析）。1 件あたりの実験要求回数（exp）とコード実行回数（code）、提出した反復の中央値（submit@）、摂動を伴わない実験要求の割合（pert0）。}
\label{tab:behavior}
{\tiny\setlength{\tabcolsep}{2.5pt}
\begin{tabular}{llrrrrrrrrrrrrrrrr}
\hline
 & & \multicolumn{4}{c}{上限 5} & \multicolumn{4}{c}{上限 10} & \multicolumn{4}{c}{上限 20} & \multicolumn{4}{c}{上限 40} \\
モデル & split & exp & code & sub@ & pert0 & exp & code & sub@ & pert0 & exp & code & sub@ & pert0 & exp & code & sub@ & pert0 \\
\hline
DeepSeek & small & \unverified{1.8} & \unverified{3.9} & \unverified{6} & \unverified{.59} & \unverified{3.3} & \unverified{7.0} & \unverified{10} & \unverified{.33} & \unverified{5.0} & \unverified{13.8} & \unverified{20} & \unverified{.26} & \unverified{7.8} & \unverified{24.6} & \unverified{39} & \unverified{.22} \\
DeepSeek & large & \unverified{1.6} & \unverified{4.0} & \unverified{6} & \unverified{.63} & \unverified{2.9} & \unverified{7.5} & \unverified{10} & \unverified{.39} & \unverified{5.9} & \unverified{14.1} & \unverified{20} & \unverified{.23} & \unverified{9.3} & \unverified{28.3} & \unverified{40} & \unverified{.21} \\
Qwen & small & \unverified{1.3} & \unverified{3.6} & \unverified{5} & \unverified{.72} & \unverified{1.7} & \unverified{8.0} & \unverified{10} & \unverified{.56} & \unverified{2.3} & \unverified{14.2} & \unverified{20} & \unverified{.40} & \unverified{2.7} & \unverified{25.6} & \unverified{36} & \unverified{.35} \\
Qwen & large & \unverified{1.3} & \unverified{3.8} & \unverified{5} & \unverified{.76} & \unverified{1.7} & \unverified{8.1} & \unverified{10} & \unverified{.59} & \unverified{2.4} & \unverified{16.3} & \unverified{20} & \unverified{.40} & \unverified{3.9} & \unverified{31.2} & \unverified{39} & \unverified{.24} \\
Nemotron & small & \unverified{2.8} & \unverified{2.9} & \unverified{6} & \unverified{.36} & \unverified{3.8} & \unverified{6.3} & \unverified{10} & \unverified{.29} & \unverified{4.7} & \unverified{12.3} & \unverified{19} & \unverified{.24} & \unverified{5.2} & \unverified{17.8} & \unverified{21} & \unverified{.26} \\
Nemotron & large & \unverified{2.7} & \unverified{3.0} & \unverified{6} & \unverified{.39} & \unverified{4.2} & \unverified{5.8} & \unverified{10} & \unverified{.27} & \unverified{5.5} & \unverified{12.2} & \unverified{19} & \unverified{.22} & \unverified{6.0} & \unverified{19.8} & \unverified{25} & \unverified{.20} \\
\hline
\end{tabular}}
\end{table}

\paragraph{平均の改善は一部の件が作る.}
表~\ref{tab:paired} に、隣接する上限で同じ件の STE がどう動いたか（両方で採点された件）を示す。最も大きい改善は 1 から 5 の段で（small の Qwen で平均 \unverified{-0.250}、DeepSeek で \unverified{-0.136}）、10 以降はどの段でも中央値がほぼ 0 で、改善した件と悪化した件が 4:3 程度に拮抗する。平均が下がるのは少数の件が大きく改善するためで、予算増は「全体が少しずつ良くなる」のではなく「解けなかった件がいくつか解ける」形で効いている。large の DeepSeek は 5 から 20 まで完全に横ばい（10 から 20 で \unverified{+0.006}）だったあと、20 から 40 で \unverified{-0.099}（改善した件 \unverified{58}\%）と最大の改善が出た。Nemotron は 10 から 20 の段で最も改善し（small \unverified{-0.086}、large \unverified{-0.071}）、40 から 80 では small \unverified{-0.026}、large \unverified{-0.016} にとどまる。

\begin{table}[t]
\centering
\caption{同じ件の対応づけ（記録外の補助分析）。隣接する上限の間で、両方で採点された件の STE の差（上限が大きい方 $-$ 小さい方）の平均と、0.01 以上改善した件・悪化した件の割合。}
\label{tab:paired}
{\tiny\setlength{\tabcolsep}{2.5pt}
\begin{tabular}{llrrrrrrrrrrrrrrr}
\hline
 & & \multicolumn{3}{c}{1$\to$5} & \multicolumn{3}{c}{5$\to$10} & \multicolumn{3}{c}{10$\to$20} & \multicolumn{3}{c}{20$\to$40} & \multicolumn{3}{c}{40$\to$80} \\
モデル & split & 平均 & 改善 & 悪化 & 平均 & 改善 & 悪化 & 平均 & 改善 & 悪化 & 平均 & 改善 & 悪化 & 平均 & 改善 & 悪化 \\
\hline
DeepSeek & small & \unverified{-.136} & \unverified{.59} & \unverified{.27} & \unverified{-.056} & \unverified{.47} & \unverified{.31} & \unverified{-.028} & \unverified{.40} & \unverified{.35} & \unverified{-.033} & \unverified{.39} & \unverified{.28} & -- & -- & -- \\
DeepSeek & large & \unverified{-.070} & \unverified{.54} & \unverified{.33} & \unverified{-.002} & \unverified{.41} & \unverified{.49} & \unverified{+.006} & \unverified{.43} & \unverified{.45} & \unverified{-.099} & \unverified{.58} & \unverified{.33} & -- & -- & -- \\
Qwen & small & \unverified{-.250} & \unverified{.69} & \unverified{.18} & \unverified{-.059} & \unverified{.45} & \unverified{.37} & \unverified{-.064} & \unverified{.47} & \unverified{.25} & \unverified{-.055} & \unverified{.43} & \unverified{.22} & -- & -- & -- \\
Qwen & large & \unverified{-.100} & \unverified{.58} & \unverified{.21} & \unverified{-.011} & \unverified{.46} & \unverified{.39} & \unverified{-.041} & \unverified{.50} & \unverified{.39} & \unverified{-.060} & \unverified{.60} & \unverified{.29} & -- & -- & -- \\
Nemotron & small & \unverified{-.017} & \unverified{.48} & \unverified{.42} & \unverified{-.005} & \unverified{.49} & \unverified{.40} & \unverified{-.086} & \unverified{.51} & \unverified{.28} & \unverified{-.002} & \unverified{.39} & \unverified{.40} & \unverified{-.026} & \unverified{.48} & \unverified{.31} \\
Nemotron & large & \unverified{+.005} & \unverified{.48} & \unverified{.45} & \unverified{-.033} & \unverified{.55} & \unverified{.40} & \unverified{-.071} & \unverified{.58} & \unverified{.32} & \unverified{-.031} & \unverified{.53} & \unverified{.38} & \unverified{-.016} & \unverified{.44} & \unverified{.42} \\
\hline
\end{tabular}}
\end{table}

\paragraph{件の難しさは安定していてモデル間で共通.}
同じモデルで上限 1 と上限 20 の件ごとの STE の順位相関（Spearman）は \unverified{+0.20}（Qwen small）から \unverified{+0.68}（Nemotron small）で、上限 1 で難しかった件は 20 でも難しい。上限 20 でのモデル間の順位相関は small で \unverified{+0.53}〜\unverified{+0.68}、large で \unverified{+0.48}〜\unverified{+0.57} で、どのモデルでも同じ件が難しい。run の中では「使った反復が多い件ほど STE が悪い」（順位相関は上限 40 の small で DeepSeek \unverified{+0.65}、Qwen \unverified{+0.68}）。これは因果ではなく、簡単な件は早く提出し、難しい件ほど最後まで粘るという行動の反映である。

\paragraph{実行で起きたこと.}
事前登録した設計からの逸脱と、実行中に判明した事実を記す。第一に、最初の投入分は取り消して全 36 run を回し直した。きっかけは 2 つで、(a) Qwen3.8-27B の large 上限 1 で、思考だけで出力上限 131,072 トークンを使い切り本文が空のまま終わる応答が \unverified{3} 件あり、当初のクライアントがこれを文脈超過として数えていたこと、(b) 一部のサーバーの起動前に run が始まり、接続先ファイルが無いことで全件が即座に失敗した run があったことである。前者は「出力上限でも本文が空なら温度 1.0 で 2 回まで引き直し、なお空なら思考枯渇として別に数える」処理を、後者は「接続先ファイルが現れるまで待つ」処理をクライアントに加え、ユーザーの判断で全 run を 1 つのコミットで回し直した。Controller・プロンプト・採点には触れていない。引き直しを入れた後、思考枯渇に至った件は 0 である。第二に、LLM の解析コードや提出モデルの評価が ODE ソルバーの C ライブラリ内で止まり、60 分の無出力で試行を打ち切ってやり直した試行が上限とともに増え、Qwen の small 上限 40 では 137 件中 \unverified{63} 件が少なくとも 1 回打ち切られた。打ち切りは反復 0 からのやり直しなので、上限 40 以上の run の所要時間の主因になっている。第三に、重みは共有ストレージに置かず各サーバーの job 開始時に取得したが、1 つの複製元に 10 サーバーが同時に接続したため 2 サーバーが外部からの再取得に戻り、起動が 1 時間半遅れた。第四に、Nemotron-3-Ultra の small 上限 20 の STE \airasval{nemotron-3-ultra-small-i20.trajectory_smape} は、先行研究で同条件を 1 回走らせた値 0.284 とほぼ一致し、run 間の変動が margin（0.03）に対して小さいことを裏付けた。

\section{考察（暫定版 v2）}
\label{sec:discussion}
\paragraph{上限 20 までは予算が性能を決める.}
論文は反復上限を 20 に固定している \cite[s1.p3]{duan-2025-measuring}。本研究の 3 モデルはこの範囲で与えられた予算を使い切っており（図~\ref{fig:it}）、STE は上限とともに下がった（図~\ref{fig:ste}）。上限 20 は性能を予算で制限している条件だと言える。ただし補助分析が示すように、追加の予算は実験ではなくコード実行に使われ、改善は一部の件に集中する。

\paragraph{Nemotron は 20 前後で自ら止まり、飽和が支持された.}
Nemotron-3-Ultra は上限 40・80 でも 20 前後で提出し、上限 40 から 80 への改善は small \airasval{nemotron-3-ultra-small-i40.trajectory_smape} $\to$ \airasval{nemotron-3-ultra-small-i80.trajectory_smape}、large \airasval{nemotron-3-ultra-large-i40.trajectory_smape} $\to$ \airasval{nemotron-3-ultra-large-i80.trajectory_smape} と小さく、Claim 11・12 が支持された。一方で上限 20 から 80 への改善は small で $-0.028$ と margin にわずかに届かず（Claim 5 反証）、large では $-0.041$ で届いた（Claim 6 支持）。small の反証は改善が無いという意味ではなく、雑音幅と同程度の改善しか無いという意味である。予算を使わないモデルにとって、上限は性能の制約ではない。

\paragraph{DeepSeek と Qwen は 40 まで改善が続き、large では 40 で初めて効く.}
DeepSeek と Qwen は上限 40 でも平均反復数が上限の 8〜10 割で、STE は 20 から 40 で small・large とも下がった。large の DeepSeek は 5 から 20 まで横ばいで、同じ件の対応づけでも 10 から 20 の差は \unverified{+0.006} と無かったが、20 から 40 では改善した件が \unverified{58}\% に増えた。反応数の多い系では、20 回の実験と解析では機構の復元に足りず、40 で初めて足りる件があると読める。上限 80 で改善が続くか（Claim 1〜4）、40 から 80 で飽和するか（Claim 7〜10）は、残る 4 run で判定する。

\paragraph{文脈長を揃えたことの効果.}
文脈長を 1M に揃えたことで、文脈超過は 32 run 合計 \unverified{30} 件（すべて large）にとどまり、先行研究で 262,144 トークンの設定のときに large の上限 20 で 1 run あたり 17 件あった超過と比べて大きく減った。large の最大プロンプト長は上限 20 で既に 100 万トークン近くに達しており、上限 80 では 1M でも超過が増える見込みである。文脈超過は提出なしとして採点されるので、上限の効果と文脈超過を分けて読む必要がある。

\section{結論（暫定版 v2）}
\label{sec:conclusion}
open-weight 3 モデルを自前 serving し、SciGym の反復上限を 1 から 80 まで振る 36 run のうち 32 run を取り込んだ時点の暫定結果を報告した。上限 1 から 20 の範囲では 3 モデルとも予算を使い切り、STE は上限とともに改善した。改善は実験回数ではなくコード実行の増加を伴い、一部の件の大きな改善として現れる。Nemotron-3-Ultra は上限 40・80 でも 20 前後で自ら提出し、40 から 80 への改善は 0.05 未満で、飽和の主張（Claim 11・12）が支持された。上限 20 から 80 への改善は large で margin を超え（Claim 6 支持）、small では雑音幅にとどまった（Claim 5 反証）。DeepSeek と Qwen は上限 40 まで改善が続き、large では 40 で初めて効いた。残る 4 run（DeepSeek と Qwen の上限 80）が揃った版で、全 36 行の表と 12 主張の判定を載せる。

\section*{データの利用可能性}
記録は \airasrecordlink{record.json} にある。実験コード、環境の固定（Dockerfile と \texttt{uv.lock}）、サーバーの起動スクリプト、各 run の結果と provenance は同じリポジトリに含まれる。補助分析の集計スクリプトは記録外であり、各 run の \texttt{instances.tar.gz} から再計算できる。

\bibliographystyle{iclr2024_conference}
\bibliography{references}

\end{document}
